\section{Controle e funções}

\begin{frame}[fragile,t]{5. IF, ELIF e ELSE}
O que é? Estruturas condicionais permitem que o programa escolha quais comandos executar de acordo com condições. \texttt{if} testa a primeira condição, \texttt{elif} testa condições adicionais e \texttt{else} executa quando nenhuma condição anterior foi satisfeita.\par
\begin{visibleenv}<2->
\EMASCcodelabel{Estrutura do comando}
\begin{lstlisting}
if condição:
    comando
elif outra_condição:
    comando
else:
    comando
\end{lstlisting}
\end{visibleenv}
\end{frame}

\begin{frame}[fragile,t]{5. IF, ELIF e ELSE / Exemplo}
\EMASCcodelabel{Exemplo prático}
\begin{lstlisting}
nota = float(input("Digite sua nota: "))

if nota >= 7:
    print("Aprovado")
elif nota >= 5:
    print("Recuperação")
else:
    print("Reprovado")
\end{lstlisting}
\begin{visibleenv}<2->
\begin{EMASCnote}
Entenda o exemplo: Se a nota for 7 ou mais, o resultado será aprovado. Se for pelo menos 5 e menor que 7, será recuperação. Abaixo de 5, será reprovado. A ordem das condições importa.
\end{EMASCnote}
\end{visibleenv}
\end{frame}

\begin{frame}[t]{6. WHILE e FOR}
São estruturas de repetição (laços), usadas para executar um bloco de comandos mais de uma vez.\par
\begin{visibleenv}<2->
\begin{EMASCnote}
Diferença prática: use \texttt{while} quando a repetição depende de uma condição; use \texttt{for} quando quiser percorrer uma sequência ou repetir com uma quantidade definida.
\end{EMASCnote}
\end{visibleenv}
\end{frame}

\begin{frame}[fragile,t]{6.1 WHILE}
O \texttt{while} repete os comandos enquanto a condição for verdadeira.\par
\begin{columns}[T,onlytextwidth]
\begin{column}{0.46\textwidth}
\EMASCcodelabel{Estrutura do comando}
\begin{lstlisting}
while condição:
    comando
\end{lstlisting}
\end{column}
\begin{column}{0.51\textwidth}
\begin{visibleenv}<2->
\EMASCcodelabel{Exemplo prático}
\begin{lstlisting}
contador = 1

while contador <= 5:
    print(contador)
    contador += 1
\end{lstlisting}
\end{visibleenv}
\end{column}
\end{columns}
\begin{visibleenv}<3->
\begin{EMASCnote}
O contador começa em 1 e aumenta uma unidade a cada repetição. Quando chega a 6, a condição se torna falsa e o laço termina. Se a variável de controle não for atualizada corretamente, pode ocorrer um laço infinito.
\end{EMASCnote}
\end{visibleenv}
\end{frame}

\begin{frame}[fragile,t]{6.2 FOR}
O \texttt{for} percorre elementos de uma sequência ou repete comandos com base em uma sequência de valores.\par
\EMASCcodelabel{Estrutura do comando}
\begin{lstlisting}
for variavel in sequência:
    comando
\end{lstlisting}
\begin{visibleenv}<2->
\EMASCcodelabel{Exemplo prático}
\begin{lstlisting}
for numero in range(1, 6):
    print(numero)
\end{lstlisting}
\end{visibleenv}
\begin{visibleenv}<3->
\begin{EMASCnote}
\texttt{range(1, 6)} produz os valores 1, 2, 3, 4 e 5. O limite final não é incluído.
\end{EMASCnote}
\end{visibleenv}
\end{frame}

\begin{frame}[fragile,t]{7. Funções}
O que é? Funções são blocos reutilizáveis de código que realizam tarefas específicas. Elas ajudam a organizar o programa e evitar repetição.\par
\begin{columns}[T,onlytextwidth]
\begin{column}{0.46\textwidth}
\EMASCcodelabel{Estrutura do comando}
\begin{lstlisting}
def nome_da_funcao():
    comandos

nome_da_funcao()
\end{lstlisting}
\end{column}
\begin{column}{0.51\textwidth}
\begin{visibleenv}<2->
\EMASCcodelabel{Exemplo prático}
\begin{lstlisting}
def saudacao():
    print("Olá! Seja bem-vindo!")

saudacao()
\end{lstlisting}
\end{visibleenv}
\end{column}
\end{columns}
\begin{visibleenv}<3->
\begin{EMASCnote}
Entenda o exemplo: A palavra \texttt{def} define a função. A linha \texttt{saudacao()} chama a função, fazendo seu bloco ser executado.
\end{EMASCnote}
\end{visibleenv}
\end{frame}

\begin{frame}[fragile,t]{7. Funções com parâmetros e retorno}
\begin{columns}[T,onlytextwidth]
\begin{column}{0.46\textwidth}
\EMASCcodelabel{Estrutura do comando}
\begin{lstlisting}
def nome_da_funcao(parametro):
    return resultado
\end{lstlisting}
\end{column}
\begin{column}{0.51\textwidth}
\begin{visibleenv}<2->
\EMASCcodelabel{Exemplo prático}
\begin{lstlisting}
def somar(a, b):
    resultado = a + b
    return resultado

total = somar(10, 5)
print(total) # 15
\end{lstlisting}
\end{visibleenv}
\end{column}
\end{columns}
\begin{visibleenv}<3->
\begin{EMASCnote}
Os parâmetros a e b recebem os valores enviados na chamada. \texttt{return} devolve o resultado para quem chamou a função.
\end{EMASCnote}
\end{visibleenv}
\end{frame}
